% Part: proof-theory
% Chapter: cut-elimination
% Section: midsequent

\documentclass[../../../include/open-logic-section]{subfiles}

\begin{document}

\olfileid{pt}{cut}{mid}

\olsection{قضیهٔ دنبالهٔ میانی و قضیهٔ هربراند}

یکی از پیامدهای مهم قضیهٔ حذف برش، قضیهٔ دنبالهٔ میانی است. این قضیه
می‌گوید اگر !!a{proof}~$\pi$ برای~$\Gamma \Sequent \Delta$ وجود داشته
باشد که در آن هر !!{formula} در صورت پیشوندی است، آنگاه !!a{proof}
~$\pi'$ای وجود دارد که شامل دنبالهٔ
$S \ident \Pi \Sequent \Lambda$، موسوم به \emph{دنبالهٔ میانی}، است؛
به‌طوری‌که همهٔ استنتاج‌های بالای~$S$ در~$\pi'$ گزاره‌ای و همهٔ
استنتاج‌های پایین~$S$ سوردارند. (!!^a{formula}~$!A$
\emph{پیشوندی}، یا در \emph{صورت پیشوندی}، است اگر به شکل
\[
\mathsf{Q}_1x_1\dots\mathsf{Q}_nx_n\,!B(x_1,\dots,x_n)
\]
باشد، که هر $\mathsf{Q}_i$ یا $\lforall$ است یا~$\lexists$.)
یکی از پیامدهای مهم قضیهٔ دنبالهٔ میانی، قضیهٔ هربراند است. بیان حالت
کلی آن کمی دشوار است، اما در کلی‌ترین صورت می‌گوید: هر مورد
!!{provable}ی که !!{sentence}~$!A$ از منطق مرتبهٔ اول باشد، همان‌گویی
گزاره‌ای~$!A'$ای وجود دارد که $!A$ را می‌توان از آن !!{prove} کرد.
صورت زیر مربوط به !!{sentence}هایی به شکل~$\lexists[x][!B(x)]$ است که
$!B$ بدون سور است:

\begin{thm}[قضیهٔ هربراند]
فرض کنید $!B(x)$ بدون سور باشد و
$\Proves \lexists[x][!B(x)]$. آنگاه ترم‌های
~$t_1$، \dots، $t_n$ای وجود دارند به‌طوری‌که
$\Proves !B(t_1) \lor \dots \lor !B(t_n)$.
\end{thm}

!!{formula}~$!B(t_1) \lor \dots \lor !B(t_n)$ را
\emph{فصل هربراند} برای~$\lexists[x][!B(x)]$ می‌نامند. چون $!B$ بدون
سور است، این فصل نیز بدون سور است. پس اگر !!{provable} باشد، با یک
برهان بدون برش !!{provable} است و در نتیجه !!a{proof}ی دارد که فقط از
استنتاج‌های گزاره‌ای استفاده می‌کند. بنابراین یک همان‌گویی است.

بخش دشوار قضیهٔ هربراند نشان‌دادن این است که برای هر مورد !!{provable}
از !!{formula}~$\lexists[x][!B(x)]$، فصلی از هربراند وجود
دارد. عکس این ادعا، یعنی اینکه اگر $\lexists[x][!B(x)]$ فصل هربراندی
داشته باشد آنگاه اثبات‌پذیر است، آسان است. در واقع، فصل هربراند
$\lexists[x][!B(x)]$ را نتیجه می‌دهد. برای نمونه، در حالت $n=2$ این
!!a{proof} را در~$\Log{G3c}$ داریم:
\[
\Axiom$!B(t_1) \fCenter \lexists[x][!B(x)], !B(t_1), !B(t_2)$
\Axiom$!B(t_2) \fCenter \lexists[x][!B(x)], !B(t_1), !B(t_2)$
\RightLabel{\LeftR{\lor}}
\BinaryInf$!B(t_1) \lor !B(t_2) \fCenter
  \lexists[x][!B(x)], !B(t_1), !B(t_2)$
\RightLabel{\RightR{\lexists}}
\UnaryInf$!B(t_1) \lor !B(t_2) \fCenter
  \lexists[x][!B(x)], !B(t_2)$
\RightLabel{\RightR{\lexists}}
\UnaryInf$!B(t_1) \lor !B(t_2) \fCenter \lexists[x][!B(x)]$
\DisplayProof
\]

برای استخراج فصل هربراند از !!a{proof}~$\pi$ برای
~$\Sequent \lexists[x][!B(x)]$، نخست حالتی را در نظر بگیرید که همهٔ
استنتاج‌های $\RightR{\lexists}$ در !!{proof} بدون برش~$\pi$ در پایان
قرار دارند:
\[
\AxiomC{}
\RightLabel{$\pi_1$}
\Deduce$ \fCenter \lexists[x][!B(x)], !B(t_1), \dots, !B(t_n)$
\RightLabel{\RightR{\lexists}}
\UnaryInf$ \fCenter \lexists[x][!B(x)], !B(t_2), \dots, !B(t_n)$
\RightLabel{$(n-1) \times \RightR{\lexists}$}
\Deduce$ \fCenter \lexists[x][!B(x)]$
\DisplayProof
\]
اگر این‌ها همهٔ استنتاج‌های $\RightR{\lexists}$ در~$\pi$ باشند که
$\lexists[x][!B(x)]$ در آن‌ها فعال است، هیچ استنتاج سوردار دیگری در
~$\pi_1$ وجود ندارد. علت این است که $!B(x)$ هیچ سور دیگری ندارد و
دنبالهٔ پایانی در سمت چپ هیچ !!{formula}ی ندارد. به بیان دیگر،
\[
\fCenter \lexists[x][!B(x)], !B(t_1), \dots, !B(t_n)
\]
دنبالهٔ میانی~$\pi$ است. افزون بر این، چون بالای دنبالهٔ میانی هیچ
استنتاج سورداری نیست، همهٔ دنباله‌های بالای آن
$\lexists[x][!B(x)]$ را در بر دارند و این فرمول نه فعال است و نه
اصلی. پس می‌توان آن را از !!{proof}~$\pi_1$ حذف کرد و !!a{proof}ی برای
~$\fCenter !B(t_1), \dots, !B(t_n)$ به دست آورد؛ سپس با $n-1$ بار
اعمال $\RightR{\lor}$، !!a{proof}ی برای فصل هربراند
$\Sequent !B(t_1) \lor \dots \lor !B(t_n)$ به دست می‌آید.

مشکل این است که نمی‌توانیم فرض کنیم حتی یک !!{proof} بدون برش برای
$\Sequent \lexists[x][!B(x)]$ همهٔ استنتاج‌های
\RightR{\lexists} خود را در یک بلوک پایانی دارد. ممکن است میان این
استنتاج‌ها، استنتاج‌هایی باشد که در آن‌ها یک $!B(t_i)$ اصلی است. از
همین رو به قضیهٔ دنبالهٔ میانی نیاز داریم.

\begin{thm}[قضیهٔ دنبالهٔ میانی]\ollabel{thm:midsequent}
اگر یک \Log{G3c}-!!{proof}~$\pi$ برای $\Gamma \Sequent \Delta$ وجود
داشته باشد که همهٔ !!{formula}های $\Gamma$ و $\Delta$ پیشوندی‌اند،
آنگاه !!a{proof}~$\pi'$ای وجود دارد که شامل دنبالهٔ
$\Pi \Sequent \Lambda$ است، به‌طوری‌که همهٔ استنتاج‌های پایین آن
سوردار و همهٔ استنتاج‌های بالای آن گزاره‌ای‌اند.
\end{thm}

\begin{proof}
\emph{مرتبهٔ} یک استنتاج سوردار در~$\pi$ را شمار استنتاج‌های گزاره‌ای
زیر آن تعریف کنید، و مرتبهٔ~$o(\pi)$ را مجموع مرتبه‌های همهٔ
استنتاج‌های سوردار آن بگیرید. اگر $o(\pi)=0$، هیچ استنتاج سورداری یک
استنتاج گزاره‌ای زیر خود ندارد و کار تمام است: چون همهٔ استنتاج‌های
سوردار یک مقدمه دارند، یک بالاترین استنتاج سوردار یکتا وجود دارد.
مقدمهٔ آن دنبالهٔ میانی است.

اگر $o(\pi)>0$، پایین‌ترین استنتاج سوردار با مرتبهٔ~$>0$ را
برگزینید. چون مرتبهٔ آن مثبت است، باید استنتاجی گزاره‌ای زیر آن باشد.
چون این استنتاج پایین‌ترینِ نوع خود است، نتیجه‌اش نمی‌تواند مقدمهٔ یک
استنتاج سوردار باشد؛ وگرنه استنتاج سوردار دوم و پایین‌تر نیز یک
استنتاج گزاره‌ای زیر خود داشت. برحسب اینکه استنتاج سوردار کدام است و
استنتاج گزاره‌ای زیر آن چیست، حالت‌های بسیار را جدا می‌کنیم. چون
دنبالهٔ پایانی فقط از فرمول‌های پیشوندی تشکیل شده است، فرمول اصلی
استنتاج سوردار نمی‌تواند در استنتاج گزاره‌ای بعدی فعال باشد. وگرنه
فرمول اصلی آن استنتاج، سوری در دامنهٔ یک رابط گزاره‌ای داشت. بنا بر
ویژگی زیر-!!{formula}، چنین فرمولی باید زیر-!!{formula} یک
!!a{formula} در دنبالهٔ پایانی می‌بود. پس !!{formula} اصلی استنتاج
سوردار !!a{element}ی از بافت استنتاج گزاره‌ای پایین‌تر است.

دو نمونه را بررسی کنیم.

نخست فرض کنید یک $\RightR{\lforall}$ در مقدمهٔ راست
\LeftR{\lif} پایان می‌یابد. آنگاه زیر-!!{proof} مربوط از~$\pi$ به
زیر-!!{proof}~$\pi_1$ زیر پایان می‌یابد:
\[
\AxiomC{}
\RightLabel{$\pi_2$}
\Deduce$\Gamma' \fCenter \Delta', \lforall[x][!A(x)], !B$
\AxiomC{}
\RightLabel{$\pi_3$}
\Deduce$!C, \Gamma' \fCenter \Delta', !A(c)$
\RightLabel{\RightR{\lforall}}
\UnaryInf$!C, \Gamma' \fCenter \Delta', \lforall[x][!A(x)]$
\RightLabel{\LeftR{\lif}}
\BinaryInf$!B \lif !C, \Gamma' \fCenter
  \Delta', \lforall[x][!A(x)]$
\DisplayProof
\]
می‌توانیم فرض کنیم $\pi$ منظم است؛ پس متغیر ویژهٔ~$c$ بیرون
~$\pi_3$ ظاهر نمی‌شود و به‌ویژه در $!B \lif !C$ یا~$\pi_2$ نیست.
اکنون برهان~$\pi_1'$ را در نظر بگیرید:
\[
\AxiomC{}
\RightLabel{$\pi_2'$}
\Deduce$\Gamma' \fCenter
  \Delta', \lforall[x][!A(x)], !A(c), !B$
\AxiomC{}
\RightLabel{$\pi_3$}
\Deduce$!C, \Gamma' \fCenter
  \Delta', \lforall[x][!A(x)], !A(c)$
\RightLabel{\LeftR{\lif}}
\BinaryInf$!B \lif !C, \Gamma' \fCenter
  \Delta', \lforall[x][!A(x)], !A(c)$
\RightLabel{\RightR{\lforall}}
\UnaryInf$!B \lif !C, \Gamma' \fCenter
  \Delta', \lforall[x][!A(x)], \lforall[x][!A(x)]$
\DisplayProof
\]
در اینجا $\pi_2'$ از تضعیف~$\pi_2$ با~$!A(c)$ در سمت راست به دست
می‌آید. (این کار هیچ استنتاجی نمی‌افزاید، به‌ویژه هیچ استنتاج سورداری
که ممکن است بر مرتبه اثر بگذارد. همچنین هیچ شرط متغیر ویژه‌ای را در
~$\pi_2$ نقض نمی‌کند، زیرا هیچ ثابتی در~$!A(c)$ به‌عنوان متغیر ویژه
در~$\pi_2$ به کار نرفته است.) شرط متغیر ویژهٔ
\RightR{\lforall} همچنان برقرار است، زیرا $c$ در~$!B \lif !C$ ظاهر
نمی‌شود.

با استفاده از پذیرفتنی‌بودن انقباض، !!a{proof}~$\pi_1''$ برای
$!B \lif !C, \Gamma' \fCenter \Delta', \lforall[x][!A(x)]$ به دست
می‌آوریم. اثبات پذیرفتنی‌بودن
$\RightR{\Contraction}$ برای $\lforall[x][!A(x)]$ و اثبات وارون‌پذیری
\RightR{\lforall} که در آن اثبات به کار رفت، مرتبهٔ هیچ استنتاجی را
تغییر ندادند؛ در بیشترین حالت فقط استنتاج‌هایی را حذف کردند. بنابراین
شمار استنتاج‌های سوردار در~$\pi_1''$ از~$\pi_1$ بیشتر نیست و هر
استنتاج سوردار در~$\pi_1''$ همان شمار استنتاج گزاره‌ای را زیر خود دارد
که استنتاج متناظر در~$\pi_1$ داشت، جز آخرین \RightR{\lforall}: در
~$\pi$ یک \LeftR{\lif} زیر آن بود که اکنون در~$\pi_1''$ به بالای آن
منتقل شده است. $\pi_1$ را در~$\pi$ با~$\pi_1''$ جایگزین کنید.
!!{proof} حاصل مرتبهٔ کمتری دارد، زیرا دست‌کم مرتبهٔ استنتاج
$\RightR{\lforall}$ حداقل~$1$ واحد کاهش یافته است. اکنون فرض استقرا
را اعمال کنید.

حالت‌هایی که استنتاج سوردار \RightR{\lexists} است حتی آسان‌ترند، زیرا
نه به انقباض نیاز است و نه نگرانی دربارهٔ شرط متغیر ویژه وجود دارد.
برای نمونه، فرض کنید پس از $\RightR{\lexists}$ یک
$\RightR{\land}$ بیاید و استنتاج نخست این بار در مقدمهٔ چپ باشد.
زیر-!!{proof} مربوط~$\pi_1$ است:
\[
\AxiomC{}
\RightLabel{$\pi_2$}
\Deduce$\Gamma' \fCenter
  \Delta', \lexists[x][!A(x)], !A(t), !B$
\RightLabel{\RightR{\lexists}}
\UnaryInf$\Gamma' \fCenter \Delta', \lexists[x][!A(x)], !B$
\AxiomC{}
\RightLabel{$\pi_3$}
\Deduce$\Gamma' \fCenter \Delta', \lexists[x][!A(x)], !C$
\RightLabel{\RightR{\land}}
\BinaryInf$\Gamma' \fCenter
  \Delta', \lexists[x][!A(x)], !B \land !C$
\DisplayProof
\]
$\pi_1$ را در~$\pi$ با~$\pi_1'$ زیر جایگزین کنید:
\[
\AxiomC{}
\RightLabel{$\pi_2$}
\Deduce$\Gamma' \fCenter
  \Delta', \lexists[x][!A(x)], !A(t), !B$
\AxiomC{}
\RightLabel{$\pi_3'$}
\Deduce$\Gamma' \fCenter
  \Delta', \lexists[x][!A(x)], !A(t), !C$
\RightLabel{\RightR{\land}}
\BinaryInf$\Gamma' \fCenter
  \Delta', \lexists[x][!A(x)], !A(t), !B \land !C$
\RightLabel{\RightR{\lexists}}
\UnaryInf$\Gamma' \fCenter
  \Delta', \lexists[x][!A(x)], !B \land !C$
\DisplayProof
\]
در اینجا $\pi_3'$ از تضعیف~$\pi_3$ با~$!A(t)$ در سمت راست به دست
می‌آید. این تضعیف فقط $!A(t)$ را به سمت راست هر دنباله در~$\pi_3$
می‌افزاید و در نتیجه مرتبهٔ هیچ استنتاج سورداری را تغییر نمی‌دهد.
مرتبهٔ استنتاج‌های سوردار در~$\pi'$ همان مرتبه در~$\pi$ است، جز
مرتبهٔ استنتاج \RightR{\lexists} که با~\RightR{\land} جابه‌جا شد و
یک واحد کاهش یافت. فرض استقرا اعمال می‌شود.
\end{proof}

\begin{prob}
حالت‌هایی از اثبات \olref{thm:midsequent} را شرح دهید که
\LeftR{\lexists} در مقدمهٔ~\RightR{\lif} پایان می‌یابد، و
\LeftR{\lforall} در مقدمهٔ چپ~\LeftR{\lor} پایان می‌یابد.
\end{prob}

\begin{proof}[اثبات قضیهٔ هربراند]
فرض کنید $\pi$ به $\Sequent \lexists[x][!B(x)]$ پایان می‌یابد. بنا بر
\olref{thm:midsequent}، می‌توانیم $\pi$ را به !!a{proof}~$\pi'$ با
دنبالهٔ میانی~$S$ تبدیل کنیم که باید به صورت زیر باشد:
\[
\Sequent \lexists[x][!B(x)], !B(t_1), \dots, !B(t_n).
\]
چون همهٔ استنتاج‌های بالای~$S$ گزاره‌ای‌اند،
$\lexists[x][!B(x)]$ نمی‌تواند در استنتاجی بالای~$S$ اصلی باشد. چون
اتمی نیست، نمی‌تواند در یک اصل اصلی باشد. چون
$\lexists[x][!B(x)]$ نمی‌تواند زیر‌فرمول سرهٔ~$!B(t_1)$ باشد
(به یاد آورید $!B(x)$ بدون سور است)، بالای~$S$ نیز نمی‌تواند فعال
باشد. بنابراین حذف $\lexists[x][!B(x)]$ از زیر-!!{proof} منتهی
به~$S$، !!{proof} درستی برای
$\Sequent !B(t_1), \dots, !B(t_n)$ به دست می‌دهد، و از آن با $n-1$
استنتاج $\RightR{\lor}$ می‌توان !!a{proof}ی برای
$\Sequent !B(t_1) \lor \dots \lor !B(t_n)$ به دست آورد.
\end{proof}

% \begin{prob}
% Find a Herbrand disjunction of $\lexists[x][\lforall[y][(!A(x) \lif
% !A(y))]]$ by finding a cut-free !!{proof} of $\Sequent
% \lexists[x][\lforall[y][(!A(x) \lif !A(y))]]$ and applying the
% transformation in \olref{thm:midsequent} to find a midsequent (if
% necessary).
% \end{prob}

\end{document}
